\documentclass{article}
\usepackage{graphicx}
\usepackage[style=ieee]{biblatex} % Establecer el estilo de las referencias como IEEE
\usepackage{xcolor}
\usepackage{hyperref}
\usepackage{titletoc}
\usepackage{adjustbox}

\hypersetup{
    colorlinks=true,
    linkcolor=blue, % Color del texto del enlace
    urlcolor=blue % Color del enlace
}

\usepackage{longtable} % Agrega el paquete longtable

\definecolor{mygreen}{RGB}{0,128,0}

\usepackage{array} % Para personalizar la tabla
\usepackage{booktabs} % Para líneas horizontales de mejor calidad
\usepackage{graphicx} % Paquete para incluir imágenes
\usepackage{float}

% Definir márgenes
\usepackage[margin=1in]{geometry}

\renewcommand{\contentsname}{\textcolor{mygreen}{Tabla de Contenidos}}

\begin{document}

\begin{titlepage}
  \centering
  % Logo de la Universidad
  \includegraphics[width=0.48\textwidth]{logo_universidad.png}
  \par\vspace{2cm}

  % Nombre de la Universidad y detalles del curso
  {\Large \textbf{Universidad Nacional de Colombia} \par}
  \vspace{0.5cm}
  {\large Ingeniería de Sistemas y Computación \par}
  {\large 2025966 Lenguajes de Programación (02)\par}
  \vspace{3cm}

  % Detalles del laboratorio y actividad
  {\large \textbf{Tarea 15} \par}
  {\large Emulación del computador diseñado\par}
  \vspace{3cm}

  % Lista de integrantes
  {\large \textbf{Integrantes:} \par}
  \vspace{0.5cm}
  \begin{tabular}{ll}
    Javier Andrés Tarazona Jiménez & jtarazonaj@unal.edu.co \\
    Eder  José Hernández Buelvas   & ehernandezbu@unal.edu.co \\
    Juan Sebastián Muñoz Lemus     & jumunozle@unal.edu.co   \\
    David Felipe Marin Rosas       & dmarinro@unal.edu.co   \\
  \end{tabular}
  \par\vspace{3cm}

  % Fecha
  {\large Mayo 7 de 2025 \par}
\end{titlepage}

\tableofcontents % Inserta la tabla de contenidos

\newpage % Salto de página para separar la tabla de contenidos del contenido del documento

% Contenido del artículo----------------------------------------------------------

%---------------------------------------------------------------------------------
% Intro --------------------------------------------------------------------------
%---------------------------------------------------------------------------------

\section{Introducción}\label{sec:intr}

El presente documento tiene como objetivo presentar el diseño e implementación de un emulador
del computador \textbf{ORISC-I}, cuyo diseño presentamos en la tarea 14. Este emulador permite simular el funcionamiento del computador, ejecutando instrucciones en un entorno controlado y visualizando el estado de los registros, memoria y otros componentes internos.

Como mencionamos en la tarea 14, el computador ORISC-I es un computador de arquitectura Von Neumann, lo que quiere decir que la memoria rom es compartida para almacenar tanto instrucciones como datos, además el computador cuenta con una serie de registros ram que se usan para almacenar datos temporales durante la ejecución de las instrucciones. El emulador también implementa el ciclo operativo de \textbf{fetch-decode-execute}, que es el ciclo básico de operación de cualquier computador, permitiendo la ejecución secuencial de instrucciones.

%---------------------------------------------------------------------------------
% Marco Teórico ------------------------------------------------------------------
%---------------------------------------------------------------------------------

\section{Marco Teórico}\label{sec:marc}

El emulador del computador ORISC-I se basa en la arquitectura Von Neumann, que define una estructura de computador donde la memoria es un espacio único para almacenar tanto instrucciones como datos. Esta arquitectura permite que el computador ejecute un conjunto de instrucciones definidas por su \textbf{Instruction Set Architecture (ISA)}, que especifica cómo se deben codificar y ejecutar las instrucciones.

El ciclo operativo de un computador se compone de tres fases principales:
\begin{itemize}
  \item \textbf{Fetch}: Recupera la instrucción desde la memoria utilizando el contador de programa (PC).
  \item \textbf{Decode}: Decodifica la instrucción para identificar su operación y operandos.
  \item \textbf{Execute}: Ejecuta la operación especificada por la instrucción, utilizando la unidad aritmético-lógica (ALU) y los registros.
\end{itemize}

En la arquitectura Von Neumann aplicada a nuestro computador ORISC-I, se cuenta con tres tipos de buses:
\begin{itemize}
  \item \textbf{Bus de datos}: Transporta los datos entre la memoria y los registros.
  \item \textbf{Bus de direcciones}: Transporta las direcciones de memoria para acceder a los datos.
  \item \textbf{Bus de control}: Transporta señales de control que coordinan las operaciones del computador.
\end{itemize}

Cuenta con una CPU compuesta por dos modulos principales:
\begin{itemize}
  \item \textbf{Unidad de Control (UC)}: Decodifica las instrucciones y genera las señales de control necesarias para la ejecución.
  \item \textbf{Unidad Aritmético-Lógica (ALU)}: Realiza operaciones aritméticas y lógicas sobre los datos almacenados en los registros.
\end{itemize}

Cuenta con 32 registros de 64 bits cada uno en la memoria del procesador, cuyas direcciones están codificadas en 5 bits, además se dispone de 4 registros especiales, ubicados en las direcciones 0 a 3:
\begin{itemize}
  \item \textbf{PC (Program Counter)}: Contador de programa, almacena la dirección de la siguiente instrucción a ejecutar.
  \item \textbf{SP (Stack Pointer)}: Puntero de pila, apunta a la parte superior de la pila en memoria.
  \item \textbf{IR (Instruction Register)}: Registro de instrucción, almacena la instrucción actualmente en ejecución.
  \item \textbf{ESTADO}: Registro de estado, contiene indicadores de condición (C, N, P, D).
\end{itemize}
Donde el resto de los registros, de la dirección 4 a la 31, son registros de propósito general.

La memoria RAM a su vez es una memoria de 24 bits de dirección, lo que permite acceder a $2^{24}$ direcciones diferentes, cada una con una palabra de 64 bits. Así, las direcciones de memoria van desde 0 hasta $2^{24}-1$ o en hexadecimal desde 0x000000 a 0xFFFFFF. La memoria se organiza en segmentos:
\begin{itemize}
  \item \textbf{Segmento de código}: Almacena las instrucciones del programa.
  \item \textbf{Segmento de datos estáticos/globales}: Almacena variables globales y datos estáticos.
  \item \textbf{Segmento de pila}: Utilizado para almacenar datos temporales y gestionar llamadas a funciones.
  \item \textbf{Segmento de entrada/salida}: Utilizado para interactuar con dispositivos externos.
\end{itemize}

Las instrucciones del ORISC-I se codifican en un formato de 64 bits, donde los primeros 6 bits representan el \textbf{opcode} (código de operación) y los siguientes 58 bits se dividen en campos que especifican los operandos y otros parámetros necesarios para la ejecución de la instrucción. El emulador implementa un conjunto de instrucciones básicas, como suma, movimiento de datos, saltos condicionales, entre otras.

\subsection{Contextualización del problema}

El emulador del computador ORISC-I se desarrolla en el contexto de la asignatura de Lenguajes de Programación, con el objetivo de proporcionar a los estudiantes una herramienta práctica para comprender los conceptos fundamentales de la arquitectura de computadores y la ejecución de instrucciones. La emulación permite experimentar con el funcionamiento interno del computador sin necesidad de hardware físico, facilitando el aprendizaje y la experimentación.

%---------------------------------------------------------------------------------
% Descripción y Justificación del Problema a Resolver ----------------------------
%---------------------------------------------------------------------------------

\section{Descripción y Justificación del Problema a Resolver}\label{sec:descr}


\subsection{Objetivo Principal}


%---------------------------------------------------------------------------------
% Diseño de la solución ---------------------------------------------------------
%---------------------------------------------------------------------------------


\section{Diseño de la solución}\label{sec:dis}

Para la implementación del emulador del computador diseñado por el grupo de trabajo, se propone
utilizar el paradigma de \textbf{programación orientada a objetos}, aplicando la arquitectura de
desarrollo \textbf{Modelo-Vista-Controlador (MVC)}. Esto permite separar la lógica funcional
(modelo), la interfaz de interacción (vista) y la gestión del flujo de ejecución (controlador),
facilitando la mantenibilidad y escalabilidad del sistema.

\subsection{Modelo}

\subsubsection*{Registros}

Esta clase representa la abstracción de los registros del sistema, incluyendo registros de
propósito general y registros especiales (PC, SP, IR, ESTADO).

\textbf{Atributos:}
\begin{itemize}
  \item \texttt{memoria} (\texttt{bool[32][64]}): Matriz que representa los 32 registros
        de 64 bits. Cada fila corresponde a un registro, y cada columna a un bit.
\end{itemize}

\subsubsection*{CPU}

Clase que modela la Unidad Central de Proceso. Coordina la ejecución del ciclo
\texttt{fetch-decode-execute} y comunica las unidades internas.

\textbf{Atributos:}
\begin{itemize}
  \item \texttt{alu} (\texttt{ALU}): Objeto de la unidad aritmético-lógica.
  \item \texttt{uc} (\texttt{UC}): Objeto de la unidad de control.
  \item \texttt{registros} (\texttt{Registros}): Banco de registros.
  \item \texttt{memoria} (\texttt{Memoria}): Memoria principal.
\end{itemize}

\textbf{Métodos:}
\begin{itemize}
  \item \texttt{\_\_init\_\_(alu, uc)}: Constructor que vincula las unidades de la CPU.
  \item \texttt{fetch()}: Carga en el registro \texttt{IR} la instrucción ubicada en la
        dirección del \texttt{PC}, y actualiza el contador de programa.
  \item \texttt{decode()}: Extrae el \texttt{opcode} y los campos relevantes de
        la instrucción para preparar su ejecución.
  \item \texttt{execute()}: Llama al método correspondiente definido en la ISA
        según el \texttt{opcode} decodificado.
\end{itemize}

\subsubsection*{UC (Unidad de Control)}

Esta clase decodifica instrucciones y genera las señales necesarias para su ejecución.

\textbf{Métodos:}
\begin{itemize}
  \item \texttt{decodificar(ir: str)}: Extrae campos de la instrucción.
  \item \texttt{seleccionar\_accion(opcode: str)}: Invoca la operación correspondiente en la ISA.
\end{itemize}

\subsubsection*{ALU}

Ejecuta operaciones aritméticas y lógicas entre registros o con valores inmediatos.

\textbf{Métodos:}
\begin{itemize}
  \item \texttt{operar(op: str, op1: int, op2: int)}: Devuelve el resultado de aplicar
        la operación \texttt{op} a los operandos.
  \item \texttt{actualizar\_estado(resultado: int)}: Actualiza los bits del
        registro \texttt{ESTADO} (C, N, P, D).
\end{itemize}

\subsubsection*{Memoria}

Clase que abstrae la memoria principal, organizada en $2^{24}$ direcciones de 64 bits cada una.

\textbf{Atributos:}
\begin{itemize}
  \item \texttt{arreglo} (\texttt{bool[$2^{24}$][64]}): Arreglo de palabras de memoria.
\end{itemize}

\textbf{Métodos:}
\begin{itemize}
  \item \texttt{leer(direccion: int) → int}: Retorna la palabra almacenada en la dirección dada.
  \item \texttt{escribir(direccion: int, palabra: int)}: Guarda una palabra de 64 bits en
        la dirección dada.
  \item \texttt{check\_structure(operator: str, direccion: int)}: Verifica si la operación es válida
        según el segmento de memoria:
        \begin{itemize}
          \item 0: Código
          \item 1: Entrada/Salida
          \item 2: Datos estáticos/globales
          \item 3: Pila
        \end{itemize}
\end{itemize}

\subsubsection*{UnidadES (Entrada/Salida)}

Abstracción de la unidad de E/S. Interactúa con dispositivos externos a través de direcciones
mapeadas en memoria.

\textbf{Métodos:}
\begin{itemize}
  \item \texttt{leer\_dispositivo(direccion: int)}: Simula lectura desde un periférico.
  \item \texttt{escribir\_dispositivo(direccion: int, valor: int)}: Simula salida de datos hacia
        un periférico.
\end{itemize}

\subsubsection*{ISA (Instruction Set Architecture)}

Define el conjunto de instrucciones que puede ejecutar el procesador. Cada instrucción tiene
su propia función que es invocada por la CPU durante la fase de ejecución.

\textbf{Métodos:}
\begin{itemize}
  \item \texttt{SUMA(rd, rb)}, \texttt{MUEV(rd, rb)}, etc.
  \item \texttt{ejecutar(opcode: str, campos: dict)}: Mapea el opcode a la función correspondiente.
\end{itemize}

% -----------------------------------------
% -----------------------------------------

\subsection{Controlador}

El controlador actúa como intermediario entre la vista y el modelo.
Es responsable de recibir las acciones del usuario, traducirlas en comandos válidos y
enviarlas a las clases del modelo, así como actualizar la vista con los cambios
de estado del sistema.

\textbf{Responsabilidades:}
\begin{itemize}
  \item Capturar eventos del usuario (ejecutar paso a paso, correr todo, resetear, cargar programa).
  \item Invocar los métodos del modelo (CPU, memoria, registros, etc.).
  \item Actualizar los datos presentados en la vista según el estado actual del sistema.
\end{itemize}

\textbf{Métodos pensados:}
\begin{itemize}
  \item \texttt{cargar\_programa(ruta)}: Traduce el programa en código máquina y lo
        almacena en memoria.
  \item \texttt{ejecutar\_paso()}: Ejecuta una instrucción y actualiza la vista.
  \item \texttt{ejecutar\_todo()}: Ejecuta hasta la instrucción \texttt{HALT}.
  \item \texttt{resetear()}: Inicializa todos los componentes del modelo.
\end{itemize}

%------------------------
%------------------------

\subsection{Vista}

La vista representa la interfaz de interacción entre el usuario y el sistema.
Es implementada como una interfaz gráfica (GUI),
Su propósito es presentar visualmente el estado interno
del computador (registros, memoria, instrucciones) y permitir la carga de programas,
control de ejecución y visualización de resultados.

\textbf{Responsabilidades:}
\begin{itemize}
  \item Mostrar en tiempo real el contenido de los registros y de la memoria.
  \item Representar gráficamente en tablas las instrucciones cargadas y el puntero de programa.
  \item Permitir al usuario:
        \begin{itemize}
          \item Cargar un programa ensamblador o en binario.
          \item Iniciar, pausar y reiniciar la ejecución.
          \item Observar los resultados de entrada/salida.
        \end{itemize}
\end{itemize}

\textbf{Ejemplo de componentes de GUI:}
\begin{itemize}
  \item Tabla de registros (R0–R31).
  \item Área de memoria con resaltado por segmento (código, datos, E/S, pila).
  \item Área de código de máquina
  \item Área de código de máquina relocalizable
  \item Área de código ensamblador
  \item Área de código alto nivel
  \item Consola interactiva para entrada/salida.
\end{itemize}



%---------------------------------------------------------------------------------
% Código Fuente ---------------------------------------------------------
%---------------------------------------------------------------------------------

\section{Código Fuente}\label{sec:cod}

El código fuente completo de este modelo se encuentra adjunto en el buzón
(11 Tarazona Jimenez Javier Andres 02.zip)
y disponible en el repositorio GitHub del proyecto:

\begin{center}
  \url{URL}
\end{center}

El repositorio contiene:
\begin{itemize}
  \item A
  \item B
  \item C
\end{itemize}

%---------------------------------------------------------------------------------
% Manual Usuario ---------------------------------------------------------
%---------------------------------------------------------------------------------

\section{Manual Usuario}\label{sec:man_u}

El primer paso es descargar el archivo \texttt{11 Tarazona Jimenez Javier Andres 02.zip}.

Una vez descargado, descomprímalo y acceda a la carpeta. Dentro de ella, cree un
entorno virtual utilizando Python 3.12. o superior. Para ello, ejecute el siguiente
comando en
la terminal o línea de comandos:

\begin{itemize}
  \item En Windows:
        \begin{verbatim}
    python3.12 -m venv nombre_del_entorno
  \end{verbatim}
  \item En macOS o Linux:
        \begin{verbatim}
    python3.12 -m venv nombre_del_entorno
  \end{verbatim}
\end{itemize}

Donde \texttt{nombre\_del\_entorno} es el nombre que desea asignar a su entorno virtual.
A continuación, active el entorno virtual:

\begin{itemize}
  \item En Windows:
        \begin{verbatim}
    .\nombre_del_entorno\Scripts\activate
  \end{verbatim}
  \item En macOS o Linux:
        \begin{verbatim}
    source nombre_del_entorno/bin/activate
  \end{verbatim}
\end{itemize}

En el archivo \texttt{constants/program.py} encontrará las constantes del programa.
En ese archivo, podrá modificar los parámetros de entrada que se detallan más abajo.\\

Después de configurar los parámetros, asegúrese de tener el entorno virtual activado.
Una vez activo, puede ejecutar el archivo principal con el siguiente comando:

\begin{center}
  \begin{adjustbox}{minipage=\linewidth, center}
    \begin{verbatim}
    python main.py
  \end{verbatim}
  \end{adjustbox}
\end{center}

%---------------------------------------------------------------------------------
% Manual Técnico ---------------------------------------------------------
%---------------------------------------------------------------------------------

\section{Manual Técnico}\label{sec:man_t}


\subsection{Fases de la Simulación}


\subsection{Manejo de Datos}

\subsection{Evaluación de la Simulación}


\subsection{Conclusiones y Recomendaciones}


%---------------------------------------------------------------------------------
% Experimentación ---------------------------------------------------------
%---------------------------------------------------------------------------------

\section{Experimentación}\label{sec:exp}

\subsubsection{Escenario 1: }

\subsubsection{Escenario 2: }

\subsubsection{Escenario 3: }


\section{Referencias}
\renewcommand{\refname}{}

\begin{thebibliography}{9}

  \bibitem{ref} \label{ref:BPS} M. Bichler, S. Merting, and A. Uzunoglu,
  “Assigning Course Schedules: About Preference Elicitation, Fairness, and Truthfulness,”
  arXiv preprint arXiv:1812.02630, 2018. [En línea]. Disponible en:
  \url{https://arxiv.org/abs/1812.02630}


\end{thebibliography}

\end{document}